\section{Inspiration behind Apollo-RTOS}

The real-time operating system developed in this project, \textit{Apollo-RTOS}, draws
inspiration from both the Apollo Guidance Computer (AGC) operating system and the
modern Linux OS. While the scheduler and error handling concepts are derived from the
AGC, the naming conventions, user interface, and implementation are influenced by
Linux.

The AGC OS was designed to meet specific mission requirements, allowing for a
simplified and optimized design. In contrast, Linux is a general-purpose OS, designed
to support applications with vastly different needs, making its architecture more
complex.

The AGC operated within a well-defined hardware and software environment, where all
components could be trusted. As a result, isolating the hardware, OS, and application
code was unnecessary. Additionally, the AGC lacked memory protection features, making
kernel/user space isolation, a fundamental aspect of general-purpose operating systems
like Linux, impossible.

In Apollo-RTOS, we adopted the same high-level design principles as the AGC. Our goal
was to implement its key innovations in a modern context while ensuring the OS could
function efficiently in a resource-constrained environment. To achieve this, we made
the following design choices:

\begin{itemize}
    \item \textbf{Hardware Selection:} The \texttt{micro:bit} v1 was chosen as the
        primary target platform. Its Cortex-M0 CPU is one of the smallest in the ARM
        family and closely resembles the AGC hardware. Additionally, it includes
        pre-installed peripheral devices, which serve as useful test cases for OS
        functionality.

    \item \textbf{No Kernel/User Space Isolation:} In microcontroller programming, the
        entire application is compiled into a single binary, meaning the kernel and
        user space are inherently combined. Since the entire codebase has full access
        to the kernel internals, enforcing strict isolation would add unnecessary
        complexity.

    \item \textbf{Monolithic Kernel Design:} Inspired by Linux, Apollo-RTOS follows a
        monolithic kernel architecture due to its simpler implementation and better
        execution efficiency in a constrained environment.

    \item \textbf{Library Calls Instead of System Calls}: Traditional system calls are
        beneficial only when user and kernel space are strictly separated. Since
        Apollo-RTOS does not enforce such isolation, we opted for direct library
        function calls, significantly simplifying implementation.

    \item \textbf{Shared Memory for Inter-Process Communication (IPC):} While shared
        memory introduces risks like race conditions and synchronization issues, it
        remains the fastest and simplest IPC method, making it the preferred choice for
        Apollo-RTOS.
\end{itemize}

Additionally, we drew inspiration from the \texttt{microbian} OS for the
\texttt{micro:bit}, as developed by \citet{microbian}. Specifically, we referenced its
hardware description headers, startup code, and interrupt handler implementations in
certain parts of our design.


% \begin{enumerate}
%     \item Hardware overview (nRF51822 microcontroller, memory, I/O)
%     \item Constraints and capabilities
%     \item Comparison with AGC hardware
%         \begin{enumerate}
%             \item Processing power
%             \item Memory availability
%             \item I/O capabilities
%             \item Power requirements 
%         \end{enumerate}
%     \item Development environment and toolchain 
% \end{enumerate}
